\documentclass[11pt]{article}
\usepackage[margin=1in]{geometry}
\usepackage{amsmath, amssymb, amsthm}
\usepackage{mathtools}
\usepackage{hyperref}
\usepackage{algorithm}
\usepackage{algpseudocode}
\usepackage{graphicx}
\usepackage{float}

\newtheorem{proposition}{Proposition}
\newtheorem{definition}{Definition}

\title{Research Notes}
\author{Senem I\c{s}\i{}k}

\begin{document}
\maketitle

\section{Bayesian Stopping with a Normal--Normal Conjugate Prior}

\subsection{Setup}

For each sample $i \in \{1,\ldots,N\}$, draw a latent mean $\mu_i$ from a prior, then draw a sequence
\[
X_1, X_2, \ldots, X_n \mid \mu_i \;\overset{\text{i.i.d.}}{\sim}\; \mathcal{N}(\mu_i, \sigma^2),
\]
with $\sigma^2$ fixed and known. We adopt the conjugate prior
\[
\mu_i \sim \mathcal{N}(\mu_0, \tau_0^2).
\]
At each time $t \in \{1,\ldots,n\}$ the agent observes $X_t$ and chooses to \emph{stop} (collect $X_t$ as terminal reward) or \emph{continue}. At $t = n$ the agent must stop.

\subsection{Three regimes}
\label{sec:regimes}

The dynamics depend on a single dimensionless ratio
\[
\rho \;:=\; \frac{\sigma^2}{\tau_0^2},
\]
the noise-to-prior precision ratio. To see why $\rho$ controls the regime, apply Proposition~\ref{prop:posterior} at $t = 1$:
\[
\mu_1
\;=\; \tau_1^2 \!\left(\frac{\mu_0}{\tau_0^2} + \frac{X_1}{\sigma^2}\right)
\;=\; \frac{\sigma^2}{\sigma^2 + \tau_0^2}\,\mu_0 \;+\; \frac{\tau_0^2}{\sigma^2 + \tau_0^2}\,X_1
\;=\; \frac{\rho}{1+\rho}\,\mu_0 \;+\; \frac{1}{1+\rho}\,X_1,
\]
so the data weight in the posterior mean is exactly $\frac{1}{1+\rho}$ after one observation, and grows monotonically toward $1$ with $t$.
\begin{enumerate}
    \item In the \textbf{diffuse-prior} regime ($\rho \ll 1$, instantiated as $\mathcal{D}_1$) the posterior is dominated by the data even after a single observation ($\frac{1}{1+\rho} \to 1$), and the plug-in baseline is asymptotically Bayes.
    \item In the \textbf{balanced} regime ($\rho \approx 1$, instantiated as $\mathcal{D}_2$) the data and prior carry equal weight in the posterior mean, neither closed-form baseline is close to optimal, and the Bayes-optimal threshold genuinely depends on the full posterior mean.
     \item In the \textbf{informative-prior} regime ($\rho \gg 1$, instantiated as $\mathcal{D}_3$) the posterior is dominated by the prior ($\frac{1}{1+\rho} \to 0$), and the prior-only baseline is asymptotically Bayes.
\end{enumerate}
We instantiate one dataset per regime in Section~\ref{sec:dataset}.

\subsection{Baselines}
\label{sec:baselines}

We compare six decision rules: four simple online policies, the Bayes-optimal online oracle (Section~\ref{sec:oracle}), and an offline hindsight upper bound. Throughout, $S_t := \sum_{s=1}^t X_s$, $\bar X_t := \frac{S_t}{t}$, and $\mu_t$ denotes the posterior mean (Proposition~\ref{prop:posterior}). The plug-in and prior-only baselines share a deterministic threshold sequence $\{\eta_t\}_{t=1}^{n-1}$ from the known-$\mu$ DP, which we derive first.

\paragraph{Known-$\mu$ thresholds.} If $\mu$ were known, the offers $X_1,\ldots,X_n$ would be uninformative about $\mu$ and the optimal continuation value $c_t^\mu := \mathbb{E}[\max\{X_{t+1}, c_{t+1}^\mu\}]$ would be a constant, with $c_{n-1}^\mu = \mu$. To evaluate the recursion we need $\mathbb{E}[\max\{X, a\}]$ for $X \sim \mathcal{N}(\mu, \sigma^2)$. Standardize with $Z := \frac{X-\mu}{\sigma}$ and $z := \frac{a-\mu}{\sigma}$; splitting at the threshold,
\[
\mathbb{E}[\max\{Z, z\}] \;=\; z\,\Pr(Z<z) + \int_z^\infty u\,\varphi(u)\,du \;=\; z\,\Phi(z) + \varphi(z),
\]
where the integral is computed via $u\,\varphi(u) = -\varphi'(u)$. Defining $\psi(z) := z\,\Phi(z) + \varphi(z)$, we obtain $\mathbb{E}[\max\{X,a\}] = \mu + \sigma\,\psi\!\left(\frac{a-\mu}{\sigma}\right)$. In standardized form $\eta_t := \frac{c_t^\mu - \mu}{\sigma}$ the recursion is $\mu$-free:
\begin{equation}
\label{eq:eta}
\eta_{n-1} = 0,
\qquad
\eta_t = \psi(\eta_{t+1}),
\qquad t = n-2,\,\ldots,\,1.
\end{equation}
The known-$\mu$ optimal rule is $\mathbf{1}[X_t \ge \mu + \sigma\eta_t]$. The two baselines below plug different estimators of $\mu$ into this rule.

\paragraph{Plug-in / frequentist.} Replace $\mu$ by its MLE $\bar X_t = \frac{S_t}{t}$:
\[
\pi_t^{\text{plug-in}}(X_{1:t}) \;=\; \mathbf{1}\!\left[\,X_t \;\ge\; \bar X_t + \sigma\,\eta_t\,\right].
\]

\paragraph{Prior-only.} Replace $\mu$ by the prior mean $\mu_0$:
\[
\pi_t^{\text{prior}}(X_t) \;=\; \mathbf{1}\!\left[\,X_t \;\ge\; \mu_0 + \sigma\,\eta_t\,\right].
\]

\paragraph{Myopic (one-step lookahead).} Stop iff the current offer beats the expected next observation: $\pi_t^{\mathrm{myopic}}(X_{1:t}) = \mathbf{1}[X_t \ge \mu_t]$ (since $\mathbb{E}[X_{t+1} \mid b_t] = \mu_t$). Coincides with $\pi^\star$ at $t = n-1$, but pathologically aggressive earlier: at $\rho = 1$, $\mu_0 = 0$ the rule reduces to $X_1 \ge X_1/2 \iff X_1 \ge 0$ on the first step, so the agent stops half the time on offer one. Empirically on $\mathcal{D}_2$, $R(\pi^{\mathrm{myopic}}) \approx 0.61$ vs $R(\pi^\star) \approx 2.04$, so myopic is strictly worse than even the secretary policy in the experiments below.

\paragraph{Bayes-optimal oracle.} The reservation policy $\pi^\star$ developed next in Section~\ref{sec:oracle}.

\paragraph{Offline optimum (hindsight).} Given $X_{1:n}$ in advance, accept the largest offer: $\pi^{\mathrm{offline}}(X_{1:n}) = \arg\max_t X_t$, with payoff $\max_t X_t$. Not implementable online, but $\mathbb{E}[\max_t X_t]$ upper-bounds every online policy. The gap $\mathbb{E}[\max_t X_t] - \mathbb{E}[X_{\tau^\star}]$ is the information cost of acting online.

\paragraph{Secretary policy.} Skip the first $r := \lfloor n/e \rfloor$ observations, set $M_r := \max_{s\le r} X_s$, and accept the first $X_t$ ($t > r$) with $X_t > M_r$. If no such $t$ exists, accept $X_n$. Ignores both the prior and the absolute scale of $X$.

\section{Bayes-optimal Oracle}
\label{sec:oracle}

\subsection{Closed-form Posterior}

\begin{proposition}
\label{prop:posterior}
Given $X_{1:t} = (X_1,\ldots,X_t)$ with $X_s \mid \mu \overset{\text{i.i.d.}}{\sim} \mathcal{N}(\mu, \sigma^2)$ and prior $\mu \sim \mathcal{N}(\mu_0, \tau_0^2)$,
\[
\mu \mid X_{1:t} \;\sim\; \mathcal{N}(\mu_t, \tau_t^2),
\qquad
\tau_t^2 = \left(\frac{1}{\tau_0^2} + \frac{t}{\sigma^2}\right)^{-1},
\qquad
\mu_t = \tau_t^2 \left(\frac{\mu_0}{\tau_0^2} + \frac{1}{\sigma^2}\sum_{s=1}^t X_s\right).
\]
\end{proposition}

\begin{proof}
By Bayes' rule, $p(\mu \mid X_{1:t}) \propto p(\mu)\, \prod_{s=1}^t p(X_s \mid \mu)$. Both factors are Gaussian, so the log-posterior is a quadratic in $\mu$. Expanding the squares and absorbing all $\mu$-free terms into the additive constant,
\begin{align*}
\log p(\mu \mid X_{1:t})
&= -\tfrac{1}{2\tau_0^2}(\mu-\mu_0)^2 - \tfrac{1}{2\sigma^2}\sum_{s=1}^t (X_s-\mu)^2 + \text{const} \\
&= -\tfrac{1}{2}\!\left(\tfrac{1}{\tau_0^2} + \tfrac{t}{\sigma^2}\right)\mu^2 + \left(\tfrac{\mu_0}{\tau_0^2} + \tfrac{1}{\sigma^2}\sum_{s=1}^t X_s\right)\mu + \text{const}.
\end{align*}
A density whose log is quadratic in $\mu$ with negative leading coefficient is Gaussian. We therefore identify the posterior with $\mathcal{N}(\mu_t, \tau_t^2)$, whose log-density expanded as a polynomial in $\mu$ is
\[
\log \mathcal{N}(\mu;\, \mu_t, \tau_t^2) \;=\; -\tfrac{1}{2\tau_t^2}\mu^2 + \tfrac{\mu_t}{\tau_t^2}\mu + \text{const}.
\]
Matching the coefficients of $\mu^2$ and $\mu$ yields
\(
\tfrac{1}{\tau_t^2} = \tfrac{1}{\tau_0^2} + \tfrac{t}{\sigma^2}
\)
and
\(
\tfrac{\mu_t}{\tau_t^2} = \tfrac{\mu_0}{\tau_0^2} + \tfrac{1}{\sigma^2}\sum_{s=1}^t X_s,
\)
which rearrange to the stated formulas.
\end{proof}

\paragraph{Recursive form.} The posterior admits a one-step update useful for online inference:
\[
\tfrac{1}{\tau_t^2} = \tfrac{1}{\tau_{t-1}^2} + \tfrac{1}{\sigma^2},
\qquad
\mu_t = \tau_t^2\!\left(\tfrac{\mu_{t-1}}{\tau_{t-1}^2} + \tfrac{X_t}{\sigma^2}\right).
\]
By Proposition~\ref{prop:posterior} the posterior depends on the data only through $S_t := \sum_{s=1}^t X_s$, so the belief state $b_t := (\mu_t, \tau_t^2)$ is sufficient for $\mu$, with $\tau_t^2$ a deterministic function of $t$ and $\mu_t$ a function of $(t, S_t)$.

\subsection{Backward Induction for the Bayes-optimal Policy}
\label{sec:dp}

Let $V_t(b_t, X_t)$ denote the optimal expected future reward at time $t$ given belief $b_t$ and current offer $X_t$, and define the continuation value
\[
C_t(b_t) \;:=\; \mathbb{E}_{X_{t+1} \mid b_t}\!\left[V_{t+1}\big(\mathrm{update}(b_t, X_{t+1}),\; X_{t+1}\big)\right].
\]
The agent chooses between stopping (payoff $X_t$) and continuing (value $C_t(b_t)$), giving $V_t(b_t, X_t) = \max\{X_t,\, C_t(b_t)\}$ for $t < n$, with forced acceptance $V_n(b_n, X_n) = X_n$. Substituting $V_{t+1}(b, x) = \max\{x, C_{t+1}(b)\}$ for $t \le n-2$ collapses the recursion to one dimension in $C$: since $\tau_t^2$ is determined by $t$ alone, $C_t$ is a function of the posterior mean $\mu$, and
\begin{equation}
\label{eq:Crec}
C_{n-1}(\mu) = \mu,
\qquad
C_t(\mu) \;=\; \mathbb{E}_{X \sim \mathcal{N}(\mu,\, \tau_t^2 + \sigma^2)}\!\Big[\max\!\big\{X,\; C_{t+1}(\mathrm{update}(\mu, X))\big\}\Big],
\end{equation}
where $\mathrm{update}(\mu, X) = \tau_{t+1}^2\!\left(\frac{\mu}{\tau_t^2} + \frac{X}{\sigma^2}\right)$ is the posterior mean from Proposition~\ref{prop:posterior}. The terminal $C_{n-1}(\mu) = \mu$ follows from $V_n = X_n$ and $\mathbb{E}[X_n \mid b_{n-1}] = \mu$.

The Bayes-optimal policy is the reservation rule
\[
\pi^\star_t(b_t, X_t) \;=\; \mathbf{1}\!\left[\, X_t \;\ge\; C_t(\mu_t) \,\right]
\quad\text{for } t < n,
\]
with forced acceptance at $t = n$.

\subsection{Approximate Dynamic Program}
\label{sec:adp}

Equation~\eqref{eq:Crec} has no elementary closed form for $t \le n - 3$, so we discretize the state space with $K$ grid points per stage and approximate the Gaussian expectation with $J$-point quadrature.

\paragraph{State space.} Fix $K, J \in \mathbb{N}$. The state we discretize is the posterior mean $\mu_t$ at stage $t$, whose marginal standard deviation across sequences is $s_t := \sqrt{\tau_0^2 - \tau_t^2}$ (by the law of total variance, $\tau_0^2 = \mathrm{Var}(\mu_t) + \mathbb{E}[\tau_t^2]$ with $\tau_t^2$ deterministic in $t$). At each stage $t \in \{1,\ldots,n-1\}$ we place a uniform grid of $K$ points on $I_t := [\mu_0 - 5 s_t,\, \mu_0 + 5 s_t]$:
\[
m_t^{(k)} \;=\; \mu_0 - 5 s_t + (k-1)\,h_t,
\qquad
h_t \;=\; \frac{10\,s_t}{K - 1},
\qquad k = 1,\ldots,K.
\]
giving a discrete state space of size $(n-1)K$.

\paragraph{Quadrature.} Use $J$-point Gauss--Hermite nodes/weights $\{(z_j, w_j)\}$ for $\int f(z)e^{-z^2}dz$. Under the change of variables $X = \mu + \sqrt{2(\tau_t^2+\sigma^2)}\,z$,
\[
\mathbb{E}_{X\sim\mathcal{N}(\mu,\,v_t)}[f(X)] \;\approx\; \frac{1}{\sqrt\pi}\sum_{j=1}^J w_j\, f\!\left(\mu + \sqrt{2 v_t}\,z_j\right),
\qquad v_t := \tau_t^2 + \sigma^2.
\]

\paragraph{Off-grid lookup with clipping.} For $\mu' \in I_t$, $\widehat{C}_t^{\mathrm{lin}}(\mu')$ is the linear interpolation of $\widehat{C}_t$ between the two adjacent grid points; for $\mu' \notin I_t$, it clips to the nearest boundary value $\widehat{C}_t(m_t^{(1)})$ or $\widehat{C}_t(m_t^{(K)})$. We interpolate $\widehat{C}$ rather than $\widehat{V}(\mu, x) = \max\{x,\widehat{C}(\mu)\}$ because $\widehat{C}$ is smooth in $\mu$ while $\widehat{V}$ has a kink.

\paragraph{Algorithm.} Define the Bayesian-update coefficient $\alpha_t := \frac{\sigma^2}{\tau_t^2 + \sigma^2}$, so that $\mathrm{update}(\mu, X) = \alpha_t\,\mu + (1-\alpha_t)\,X$. Algorithm~\ref{alg:adp} runs the backward induction; it precomputes $\tau_t^2$, $h_t$, $v_t$, $\alpha_t$ for all $t$, the grids $\{m_t^{(k)}\}$, and the Gauss--Hermite nodes/weights $\{(z_j, w_j)\}_{j=1}^J$.

\begin{algorithm}[H]
\caption{Approximate DP for the Bayes-optimal reservation price.}
\label{alg:adp}
\begin{algorithmic}[1]
\State $\widehat{C}_{n-1}(m_{n-1}^{(k)}) \gets m_{n-1}^{(k)}$ \quad for $k = 1,\ldots,K$
\For{$t = n-2,\, n-3,\, \ldots,\, 1$}
  \For{$k = 1, \ldots, K$}
    \State $S \gets 0$
    \For{$j = 1, \ldots, J$}
      \State $x_j \gets m_t^{(k)} + \sqrt{2 v_t}\,z_j$
      \State $\mu_j' \gets \alpha_t\, m_t^{(k)} + (1-\alpha_t)\, x_j$
      \State $S \gets S + w_j\,\max\!\big\{x_j,\; \widehat{C}_{t+1}^{\mathrm{lin}}(\mu_j')\big\}$
    \EndFor
    \State $\widehat{C}_t(m_t^{(k)}) \gets S / \sqrt{\pi}$
  \EndFor
\EndFor
\State \Return $\{\widehat{C}_t(m_t^{(k)})\}_{t,k}$
\end{algorithmic}
\end{algorithm}

\noindent The deployed policy at inference is $\widehat\pi^\star_t(b_t, X_t) = \mathbf{1}[X_t \ge \widehat{C}_t^{\mathrm{lin}}(\mu_t)]$.

\paragraph{Complexity per DP timestep.} Each stage updates $K$ grid values, each costing $O(J)$, and reads/writes $O(K)$ memory: $O(KJ)$ time, $O(K)$ space per stage.

\section{Dataset and Labeling}
\label{sec:dataset}

\subsection{Three datasets}

We instantiate the three regimes of Section~\ref{sec:regimes} at $\rho \in \{0.1,\, 1,\, 10\}$, well-separated on a log scale (data weights $\frac{1}{1+\rho} \in \{0.91,\, 0.50,\, 0.09\}$ at $t = 1$):
\begin{center}
\small
\begin{tabular}{l c c l}
\hline
dataset & $\rho$ & data weight at $t=1$ & regime \\
\hline
$\mathcal{D}_1$ & $0.1$ & $0.91$ & diffuse-prior (data-dominant) \\
$\mathcal{D}_2$ & $1$   & $0.50$ & balanced \\
$\mathcal{D}_3$ & $10$  & $0.09$ & informative-prior (prior-dominant) \\
\hline
\end{tabular}
\end{center}

\subsection{Hyperparameters}

We fix $\mu_0 = 0$, $\sigma = 1$, and horizon $n = 64$, varying $\tau_0 \in \{\sqrt{10},\, 1,\, 1/\sqrt{10}\}$ to realize the three values of $\rho$. Training sequences are streamed from each generative model on the fly (a fresh batch each step); we hold out fixed validation and test sets of $N_{\mathrm{val}} = N_{\mathrm{test}} = 10^4$ sequences per dataset, sampled with distinct RNG seeds from the same generative model. We solve the approximate DP exactly once per dataset with $K = 256$ and $J = 32$ and reuse the threshold tables to label every sequence; each solve runs in well under a second and matches the $K = 1024,\, J = 64$ table to better than $10^{-2}$. The Gauss--Hermite convergence rate is capped by the kink in the integrand $\max\{x,\, \widehat{C}_{t+1}(\mathrm{update}(\mu, x))\}$ at the indifference point; refining $K$ or $J$ further hits a wall around $1.5 \times 10^{-2}$ without split-at-kink quadrature. This is small relative to the $\sim\sigma$ spread of $C_t(\mu_t)$ and well below the model's training noise, so we accept it.

\subsection{Labeling}

Given the precomputed table $\{\widehat{C}_t(m_t^{(k)})\}$, labeling a sequence $X_{1:n}$ in $\mathcal{D}_i$ proceeds by: for each $t \in \{1,\ldots,n-1\}$, update the posterior mean $\mu_t$ via the conjugate recursion of Proposition~\ref{prop:posterior}, evaluate the precomputed threshold $\widehat{C}_t^{\mathrm{lin}}(\mu_t)$ (Section~\ref{sec:adp}), and emit the two labels
\[
y_t^{\mathrm{cv}} \;:=\; \widehat{C}_t^{\mathrm{lin}}(\mu_t)
\quad\text{(continuation value)},
\qquad
y_t^{\mathrm{act}} \;:=\; \mathbf{1}\!\left[X_t \ge \widehat{C}_t^{\mathrm{lin}}(\mu_t)\right]
\quad\text{(action)}.
\]
At step $t = n$ we set $y_n^{\mathrm{act}} = 1$ (forced acceptance) and leave $y_n^{\mathrm{cv}}$ undefined. Per-sequence labeling is $O(n)$: one conjugate update and one constant-time grid lookup per step.

We retain both label streams to train and evaluate the transformer under each separately --- a sparse classification target ($y^{\mathrm{act}}$, deployment-aligned) versus a dense regression target ($y^{\mathrm{cv}}$, the full continuous threshold) --- to isolate whether dense supervision yields a better policy than sparse, a central question for the experiments below.

\section{Model}
\label{sec:model}

\subsection{Architecture}

We use a GPT-2 decoder with depth $L = 8$, embedding dimension $d_{\mathrm{emb}} = 128$, and $M = 4$ attention heads, with a causal (lower-triangular) attention mask and no dropout (training sequences are streamed, so over-fitting is not a concern). The model has roughly $1.6$M parameters.

\subsection{Input and output projections}

Each scalar observation $X_t$ is mapped to the embedding by a learnable affine layer $\mathbf{e}_t = W_{\mathrm{in}}\, X_t + \mathbf{b}_{\mathrm{in}}$ ($W_{\mathrm{in}}, \mathbf{b}_{\mathrm{in}} \in \mathbb{R}^{d_{\mathrm{emb}}}$), followed by standard learned positional embeddings.

The hidden state $\mathbf{h}_t \in \mathbb{R}^{d_{\mathrm{emb}}}$ at each position $t$ feeds two linear heads: a continuation-value head $\widehat{C}_t = \mathbf{w}_{\mathrm{cv}}^\top \mathbf{h}_t + b_{\mathrm{cv}}$ predicting the oracle threshold, and an action head $\hat a_t = \mathrm{sigmoid}(\mathbf{w}_{\mathrm{act}}^\top \mathbf{h}_t + b_{\mathrm{act}})$ predicting acceptance probability. The causal mask makes $\mathbf{h}_t$ depend only on $X_{1:t}$, matching the deployment information set.

\subsection{Training objective}

We train one model per dataset $\mathcal{D}_i$ and per supervision signal --- six models total (three continuation-value, three action), each carrying only the corresponding output head.

For a model with parameters $\theta$ producing predictions $\widehat{C}_t(X_{1:t};\theta)$ or $\hat a_t(X_{1:t};\theta)$, the per-sequence losses are
\begin{align*}
\ell_{\mathrm{cv}}(\theta;\, X_{1:n})
&\;=\; \frac{1}{n-1}\sum_{t=1}^{n-1}\Big(\widehat{C}_t(X_{1:t};\theta) \;-\; y_t^{\mathrm{cv}}\Big)^{\!2}, \\
\ell_{\mathrm{act}}(\theta;\, X_{1:n})
&\;=\; -\frac{1}{n-1}\sum_{t=1}^{n-1}\!\Big[\, y_t^{\mathrm{act}}\,\log \hat a_t(X_{1:t};\theta) \;+\; (1 - y_t^{\mathrm{act}})\,\log\!\big(1 - \hat a_t(X_{1:t};\theta)\big)\,\Big],
\end{align*}
where $y_t^{\mathrm{cv}}$ and $y_t^{\mathrm{act}}$ are the oracle labels from Section~\ref{sec:dataset}. The sums run only to $n-1$: step $n$ is forced acceptance and carries no learning signal ($y_n^{\mathrm{act}} = 1$ deterministically; $y_n^{\mathrm{cv}}$ is undefined).

We minimize the expected per-sequence loss over $\mathcal{D}_i$,
\[
\theta^\star_{\mathrm{cv}} \;=\; \arg\min_\theta\; \mathbb{E}_{X_{1:n}\sim\mathcal{D}_i}\!\left[\,\ell_{\mathrm{cv}}(\theta;\,X_{1:n})\,\right],
\qquad
\theta^\star_{\mathrm{act}} \;=\; \arg\min_\theta\; \mathbb{E}_{X_{1:n}\sim\mathcal{D}_i}\!\left[\,\ell_{\mathrm{act}}(\theta;\,X_{1:n})\,\right],
\]
estimated by mini-batches of fresh sequences. We use Adam with peak learning rate $10^{-4}$, no weight decay, batch size 64, $5\times 10^5$ training steps, learned absolute position embeddings, and a cosine learning-rate schedule with linear warmup over the first $10^5$ steps (20\% of training). Biases, weights, and position embeddings are initialized with uniform scaling. We checkpoint periodically during training and use the lowest-validation-loss checkpoint for the experiments in Section~\ref{sec:experiments}.

\section{Experiments}
\label{sec:experiments}

\subsection{Metrics}

For a policy $\pi$ on test sequences from $\mathcal{D}$, with stopping time $\tau_\pi$, the primary metric is the expected payoff
\[
R(\pi;\,\mathcal{D}) \;:=\; \mathbb{E}_{X_{1:n}\sim\mathcal{D}}\!\left[X_{\tau_\pi}\right],
\]
estimated on a held-out test split. To diagnose what the model has learned --- not merely how well it scores --- we additionally report its per-step action-agreement with each baseline $\pi'$,
\[
\mathrm{Agree}(\pi_\theta,\, \pi';\,\mathcal{D}) \;:=\; \mathbb{E}\!\left[\,\frac{1}{n-1}\sum_{t=1}^{n-1}\mathbf{1}\!\left[\pi_{\theta,t}(X_{1:t}) = \pi'_t(X_{1:t})\right]\,\right].
\]

\subsection{In-distribution evaluation}
\label{sec:exp-indist}

Each model is evaluated on its training distribution and compared against the six baselines. 

\textbf{Finding.} All six models recover the oracle policy in-distribution: $R/R^\star \in [0.99981,\, 1.00006]$, oracle agreement $\ge 0.9998$ on $10^4$ test sequences $\times\,63$ decision steps (Table~\ref{tab:indist}). Both supervision signals (cv regression, act classification) are sufficient.

\begin{table}[H]
\centering
\small
\begin{tabular}{l c c c c c c c c}
\hline
model & train & sup. & $R(\pi_\theta)$ & $R(\pi^\star)$ & gap & agree($\pi^\star$) & $R/R^\star$ & $R/R^{\mathrm{off}}$ \\
\hline
D1\_cv  & $\mathcal{D}_1$ & cv  & 2.0119 & 2.0118 & $+0.000038$ & $0.999970$ & $1.000019$ & $0.869352$ \\
D2\_cv  & $\mathcal{D}_2$ & cv  & 2.0422 & 2.0422 & $+0.000000$ & $0.999989$ & $1.000000$ & $0.873908$ \\
D3\_cv  & $\mathcal{D}_3$ & cv  & 2.0792 & 2.0792 & $+0.000009$ & $0.999989$ & $1.000004$ & $0.887060$ \\
D1\_act & $\mathcal{D}_1$ & act & 2.0115 & 2.0118 & $-0.000382$ & $0.999865$ & $0.999810$ & $0.869170$ \\
D2\_act & $\mathcal{D}_2$ & act & 2.0420 & 2.0422 & $-0.000127$ & $0.999924$ & $0.999938$ & $0.873853$ \\
D3\_act & $\mathcal{D}_3$ & act & 2.0794 & 2.0792 & $+0.000132$ & $0.999890$ & $1.000063$ & $0.887112$ \\
\hline
\end{tabular}
\caption{In-distribution test-set payoff and per-step oracle agreement, $N_{\mathrm{test}} = 10^4$. $R^{\mathrm{off}} = \mathbb{E}[\max_t X_t]$ is offline hindsight; $R/R^{\mathrm{off}} \approx 0.87 \to 0.89$ across $\mathcal{D}_1 \to \mathcal{D}_3$ is the structural cost of acting online, not a model deficiency.}
\label{tab:indist}
\end{table}

Table~\ref{tab:baselines} reports the test-split payoff of every baseline on each dataset, for context.

\begin{table}[H]
\centering
\small
\begin{tabular}{l c c c c c c}
\hline
dataset & $\pi^\star$ & $\pi^{\mathrm{plug\text{-}in}}$ & $\pi^{\mathrm{prior}}$ & $\pi^{\mathrm{myopic}}$ & $\pi^{\mathrm{secretary}}$ & $\pi^{\mathrm{offline}}$ \\
\hline
$\mathcal{D}_1$ ($\rho = 0.1$) & 2.0118 & 2.0109 & 0.6816 & 0.4043 & 1.3164 & 2.3142 \\
$\mathcal{D}_2$ ($\rho = 1$)   & 2.0422 & 2.0335 & 1.5140 & 0.6133 & 1.3390 & 2.3368 \\
$\mathcal{D}_3$ ($\rho = 10$)  & 2.0792 & 2.0406 & 2.0151 & 0.7597 & 1.3461 & 2.3440 \\
\hline
\end{tabular}
\caption{Test-split expected payoff of each baseline on each dataset, $N_{\mathrm{test}} = 10^4$. Plug-in is asymptotically optimal on $\mathcal{D}_1$ and tracks the oracle to within $0.04$ on $\mathcal{D}_3$; prior-only is asymptotically optimal on $\mathcal{D}_3$ and collapses to $0.68$ on $\mathcal{D}_1$.}
\label{tab:baselines}
\end{table}

\subsection{Out-of-distribution evaluation}
\label{sec:exp-ood}

We evaluate each model on the two regimes it was not trained on. The diagnostic question: has the model internalized the regime-invariant Bayesian algorithm, or merely the asymptotically-optimal policy of its training regime? The two are observationally indistinguishable in-distribution but pull apart OOD.

\textbf{Finding: shortcut, not algorithm.} OOD payoff drops below the oracle, and per-step agreement with the training-regime shortcut exceeds agreement with the eval-distribution oracle. cv vs act makes essentially no difference. Figures~\ref{fig:ood_heatmap}, \ref{fig:agreement_matrix}, \ref{fig:payoff_bars} show the full $2 \times 3$ matrix of $(\text{train}, \text{eval})$ pairs; the headline numbers are
\begin{itemize}
  \item \textbf{$\mathcal{D}_3 \to \mathcal{D}_1$ (largest failure).} cv $R/R^\star = 0.582$, act $0.621$. $\mathrm{agree}(\pi^{\mathrm{prior}}) \ge \mathrm{agree}(\pi^\star)$: $0.86 \ge 0.84$ (cv), $0.83 \ge 0.83$ (act) --- the $\mathcal{D}_3$-trained model has memorized prior-only and applied it where it scores $R/R^\star = 0.34$.
  \item \textbf{$\mathcal{D}_1 \to \mathcal{D}_3$ (smallest failure).} cv $R/R^\star = 0.982$, act $0.982$. $\mathrm{agree}(\pi^{\mathrm{plug}}) > \mathrm{agree}(\pi^\star)$: $0.998 > 0.992$ --- same memorization signature, milder because plug-in itself is within $R/R^\star = 0.98$ of the oracle on $\mathcal{D}_3$.
  \item \textbf{$\mathcal{D}_2$-trained, intermediate.} cv $R/R^\star = 0.892$ ($\to \mathcal{D}_1$), $0.987$ ($\to \mathcal{D}_3$); act tracks within $0.003$. The no-shortcut training regime does not protect against OOD failure --- the model learns a $\mathcal{D}_2$-shaped threshold that does not transport.
  \item \textbf{cv vs act null.} $|\Delta R/R^\star| \le 0.04$ in every cell, $\le 0.01$ in five of six. Dense supervision does not help OOD generalization.
\end{itemize}

\begin{figure}[H]
\centering
\includegraphics[width=0.95\linewidth]{fig_ood_heatmap}
\caption{OOD generalization gap $R(\pi_\theta) - R(\pi^\star)$ per (supervision, train) subplot. In-distribution cells sit at zero; the asymmetric OOD failure pattern peaks at prior-dominant $\to$ data-dominant ($-0.84$, $R/R^\star \approx 0.58$).}
\label{fig:ood_heatmap}
\end{figure}

\begin{figure}[H]
\centering
\includegraphics[width=0.95\linewidth]{fig_agreement_matrix}
\caption{Per-step action agreement of each model with bayes\_optimal, plug\_in, and prior\_only on each eval distribution. OOD cells reveal which policy the model has actually learned: for $\mathcal{D}_3$-trained models on $\mathcal{D}_1$, $\mathrm{agree}(\pi^{\mathrm{prior}}) \ge \mathrm{agree}(\pi^\star)$.}
\label{fig:agreement_matrix}
\end{figure}

\begin{figure}[H]
\centering
\includegraphics[width=0.95\linewidth]{fig_payoff_bars}
\caption{Test-split expected payoff per policy and per (train, eval) cell. First three groups are in-distribution; remaining six are OOD. Bayes-optimal sets the ceiling within each group.}
\label{fig:payoff_bars}
\end{figure}


\subsection{Threshold trajectories}
\label{sec:exp-thresh}

For threshold-form policies $\pi_t(X_{1:t}) = \mathbf{1}[X_t \ge \theta_t]$, we plot the per-step threshold (mean $\pm 1$ std over test sequences) as a function of $t$. Models contribute their cv-head output $\widehat{C}_t(X_{1:t};\theta)$; act-supervised models are excluded since they emit a binary logit, not a scalar threshold. Figure~\ref{fig:threshold_v1} shows the full baseline landscape including the offline hindsight ceiling; Figure~\ref{fig:threshold_v2} hard-zooms to $(1.5, 2.2)$ on the three threshold-form online baselines plus the cv models for legibility.

\textbf{Finding.} In-distribution, each cv model's threshold curve sits on its training-distribution oracle, well below the offline ceiling --- a proper online threshold, not a collapse onto hindsight. OOD, the curves continue to look like the training-distribution policy transplanted onto the new input range; the $\mathcal{D}_3 \to \mathcal{D}_1$ panel collapses near $\mu_0 = 0$, reproducing the §\ref{sec:exp-ood} shortcut diagnostic at the level of per-step decision boundaries.

\begin{figure}[H]
\centering
\includegraphics[width=0.95\linewidth]{threshold_trajectories}
\caption{Threshold trajectories on each eval distribution, native y-range. Threshold-form baselines (bayes\_optimal, plug\_in, prior\_only, myopic) drawn as mean $\pm 1$ std; secretary from $t > \lfloor n/e \rfloor$ onward; offline as a horizontal dashed line at $\mathbb{E}[\max_t X_t]$ (hindsight ceiling). cv models in red/cyan/olive dashed.}
\label{fig:threshold_v1}
\end{figure}

\begin{figure}[H]
\centering
\includegraphics[width=0.95\linewidth]{threshold_trajectories_v2}
\caption{Same content as Figure~\ref{fig:threshold_v1}, restricted to bayes\_optimal / plug\_in / prior\_only and the cv models, y-axis hard-zoomed to $(1.5, 2.2)$. D1\_cv red dashed, D2\_cv orange solid, D3\_cv purple dotted; out-of-band curves are clipped and listed in the off-axis annotation. D3\_cv on the data-dominant panel collapses near $\mu_0 \approx 0$ --- its disappearance below the visible band is the OOD failure signature.}
\label{fig:threshold_v2}
\end{figure}

\section{Next iteration: harder problem instances}
\label{sec:next-iteration}

The findings of Section~\ref{sec:experiments} establish a clean qualitative picture, but four artifacts of the current setup limit it. We re-run the pipeline at a harder scale to address them.

\paragraph{Compressed dynamic range.} With $\sigma = 1$ the marginal $X$ standard deviation $\sqrt{\sigma^2 + \tau_0^2}$ ranges from $1.05$ on $\mathcal{D}_3$ to $3.32$ on $\mathcal{D}_1$, and on $\mathcal{D}_3$ payoff differences live near the floating-point resolution floor (e.g., the cv-vs-act in-distribution gap is $O(10^{-4})$).

\paragraph{Asymmetric OOD detection.} The shortcut-vs-eval-oracle gap is $1.33$ in $\mathcal{D}_3 \to \mathcal{D}_1$ but only $0.04$ in $\mathcal{D}_1 \to \mathcal{D}_3$, so the cleaner direction has too small a measurement window to distinguish shortcut from algorithm.

\paragraph{Non-identifiability of $\tau_0$.} With $\tau_0$ varying across regimes and $\sigma$ fixed, the model has no in-context signal to identify the regime --- $\tau_0$ governs cross-sequence spread of $\mu$ and is invisible from a single sequence. OOD failure is therefore structurally inevitable for any model lacking out-of-band information about the regime.

\paragraph{Easy task.} Calibration showed payoff matching the oracle by step $55{,}000$ at $n = 64$, well within the $5\times 10^5$ training budget. The architecture is overprovisioned and in-distribution success is not a meaningful test.

\subsection{V2 spec}

The conjugate-prior math, the DP, and the architecture are unchanged; only the parameters change.
\begin{center}
\small
\begin{tabular}{l c c c c}
\hline
dataset & $\tau_0^2$ & $\sigma^2$ & $\sigma$ & $\rho$ \\
\hline
$\mathcal{D}_1$ (data-dominant)    & $100$ & $1$     & $1$   & $0.01$ \\
$\mathcal{D}_2$ (balanced)         & $100$ & $100$   & $10$  & $1$ \\
$\mathcal{D}_3$ (prior-dominant)   & $100$ & $10000$ & $100$ & $100$ \\
\hline
\end{tabular}
\end{center}
We fix $\tau_0 = 10$ and vary $\sigma \in \{1, 10, 100\}$, increase $n$ from $64$ to $256$, and bump the ADP resolution to $K = 1024$, $J = 64$.

The key change is varying $\sigma$ rather than $\tau_0$: the within-sequence sample variance is a consistent estimator of $\sigma^2$, so the regime variable is now \emph{identifiable from a single in-context sequence}. A model that has learned the algorithm should adapt its threshold to the in-context variance estimate; a shortcut-memorizing model cannot. The OOD diagnostic gains a directionally observable prediction. The $100\times$ separation between $\mathcal{D}_1$ and $\mathcal{D}_3$ removes the asymmetric-detection limitation, the larger horizon $n = 256$ stops in-distribution success from being too easy, and the larger marginal $X$ scales lift everything off the precision floor.

\subsection{Predictions}

In-distribution oracle-matching should still hold. OOD payoffs should drop in both directions, large enough to resolve cleanly under either hypothesis. If the algorithm hypothesis is even partially correct, the variance-identifiability of the regime variable should yield smaller OOD gaps than the v1 results suggest. The cv-vs-act null result is the most likely v1 finding to break: at $n = 256$ dense supervision may matter where it didn't at $n = 64$. Attention should develop a new specialization for the centered second moment $\frac{1}{t}\sum (X_s - \bar X_t)^2$, distinct from the uniform running-sum head; this pattern is absent from v1 but required by the algorithm at v2 scale.

If a v1 finding fails to replicate, we will report the non-replication explicitly.
\end{document}
